Let $F$ be a filtered object in $\cC$. Recall from \cref{def:SS_Fil} that we have the spectral sequence $\{ E_r^{n,p}, d_r^{n,p} \}$ associated with $F$. To understand the differentials $d_r^{n,p}$, we consider the following commutative diagram arising from the connecting morphisms in the long exact sequences of homotopy groups:

\begin{equation}\label{eq:diff_well_def}
\xymatrix{
\pi_n\left( \nicefrac{F_p}{F_{p - r}} \right) \ar[r]\ar[d]_{\delta_{p - r}} & \pi_n\left( \nicefrac{F_p}{F_{p - 1}} \right) \ar[r] & \pi_n\left( \nicefrac{F_{p + r - 1}}{F_{p - 1}} \right) \ar[d]^{\delta_{p - 1}} \\
\pi_{n - 1}\left( \nicefrac{F_{p - r}}{F_{p - 2r}} \right) \ar[r] & \pi_{n - 1}\left( \nicefrac{F_{p - r}}{F_{p - r - 1}} \right) \ar[r] & \pi_{n - 1}\left( \nicefrac{F_{p - 1}}{F_{p - r - 1}} \right),
}
\end{equation}

where the vertical maps $\delta_{p - r}$ and $\delta_{p - 1}$ are the connecting homomorphisms from the long exact sequences associated with the cofiber sequences in the filtration.

Let $x \in E_r^{n,p}$ be an element that survives to the $r$-th page, and let $x' \in \pi_n\left( \nicefrac{F_p}{F_{p - r}} \right)$ be a lift of $x$. By the commutativity of Diagram~\Cref{eq:diff_well_def}, the image of $\delta_{p - r}(x')$ in $\pi_{n - 1}\left( \nicefrac{F_{p - 1}}{F_{p - r - 1}} \right)$ depends only on $x$ and not on the specific choice of the lift $x'$. This property allows us to define the differential $d_r^{n,p}$.

\begin{defn}
Let $F$ be a filtered spectrum with spectral sequence $\rnp{E}$ and let $x \in E_r^{n,p}$. Let $x' \in \pi_n\left( \nicefrac{F_p}{F_{p - r}} \right)$ be a lift of $x$. We define the differential
\begin{equation*}
    d_r^{n,p}(x) \in E_r^{n - 1, p - r}    
\end{equation*}
to be the image of $\delta_{p - r}(x')$ in $E_r^{n - 1, p - r}$ under the natural map induced by the composition
\begin{equation*}
    \pi_{n - 1}\left( \nicefrac{F_{p - r}}{F_{p - 2r}} \right) \to  E_r^{n - 1, p - r}.    
\end{equation*}
Diagrammatically, this is represented as:
\begin{equation}\label{eq:differential_diagram}
\xymatrix{
x' \ar[d]_{\delta_{p - r}} & \ar@{<-}[l] x \\
\delta_{p - r}(x') \ar@{|->}[r] & d_r^{n,p}(x).
}
\end{equation}
\end{defn}

The definition ensures that $d_r^{n,p}(x)$ is well-defined, independent of the choice of lift $x'$.


\subsubsection*{Chain-level description of the differential}

Similar to the description of the pages, as maps between chain complexes, we can unwind Lurie's description for filtered spectrum, to get a chain complex description of the differential.

To facilitate the proof of \cref{lem:Ch_diff}, we first introduce the following definition and lemma.

\begin{defn}
    We denote the "shift to the left" functor by 
    \begin{equation*}
        T \colon \Ch(\Sp) \to \Ch(\Sp)
    \end{equation*}
    which is given by precomposing with 
    \begin{equation*}
        \operatorname{-1}: \bbZ \to \bbZ, \quad n \mapsto n-1
    \end{equation*}
\end{defn}

\begin{lem} \label{lem:shift_map}
    Let $L \in \Ch_{(\infty,n]}(\Sp)$ and $K \in \Ch_{[n,-\infty)}(\Sp)$ be two chain complexes. Then
    \begin{equation*}
        \Map(\S^{-1} T L,K) \simeq \Map(L, K)
    \end{equation*}  
\end{lem}
\begin{proof}
    First, note that $\Fil(\Sp)$ is generated (under colimits) by $\yo_p$. Passing to $\Ch(\Sp)$, we get that $\Ch(\Sp)$ is generated by $\cI \yo_p = \bbS_p^{-p}$.
    Using this observation, it suffices to show that:
    \begin{equation*}
        \Map(\bbS_{p}^{-p}, K) \simeq \Map(\S^{-1} T \bbS_{p}^{-p},K) \simeq \Map(S_{p+1}^{-(p+1)},K)
    \end{equation*}
    and by taking $\cI$, our goal is to show:
    \begin{equation*}
        \Map(\yo_p, \cI K) \simeq \Map(\yo_{p+1},\cI K)
    \end{equation*}
    which is equivalent to
    \begin{equation}\label{eq:tuncatedChainLem}
        (\cI K)_p \simeq (\cI K)_{p+1}
    \end{equation}
    Note that the assumption on $L$ translates to $p \geq n$.

    Let $F = \cI K$, where $K \in \Ch_{[n,-\infty)}(\Sp)$. Then for $q \geq n+1$ we have
    \begin{equation*}
        0 = K_q = (\cA F)_q = \S^{-q} F_q /F_{q-1}
    \end{equation*}
    which shows that \cref{eq:tuncatedChainLem} holds.
\end{proof}

\begin{prop}\label{lem:Ch_diff}
    Let $K \in \Ch(\Sp)$ be a chain complex in Spectra, let $F = \cI K$ be its associated filtered spectrum, and let $\rnp{E}$ be the spectral sequence of $F$. Let $x \in \rnp{\cE}$ be an element surviving to the $r$-th page (see \cref{def:survive}). Then the differential $\partial_r(x)$ is given by the composition of chain maps:
    \begin{equation*}
        \xymatrix
        {   \mathbb{S}^{n-p} \ar[r]\ar[d]_x & \dots \ar[r]\ar[d] &  0 \ar[d]\\
            K_p \ar[r]\ar[d] & \dots \ar[r]\ar[d] &  K_{p-r+1} \ar[d] \\
            0 \ar[r] & \dots \ar[r] & K_{p-r} 
        }
    \end{equation*}
    viewed as a map $\bbS^{n-p} \to K_{p-r}$, where the first map is an extension of $x$ (see \cref{lem:ChSur}), and the second is induced by the structure of the chain complex $K$.
\end{prop}

\begin{proof}
    We begin with Lurie's description, \cref{eq:diff_Fil}. The differential is given by the map:
    \begin{equation}
        \pi_n F_p/F_{p-r} \to \pi_{n-1} F_{p-r}/F_{p-2r}
    \end{equation}
    hence $\partial_r$ is induced by post-composing with the connecting morphism:
    \begin{equation*}
        F_p/F_{p-r} \to \Sigma F_{p-r}/F_{p-2r}
    \end{equation*}
    which comes from the cofiber sequence:
    \begin{equation*}
        F_{p-r}/F_{p-2r} \to F_p/F_{p-2r} \to F_p/F_{p-r}.
    \end{equation*}

    As we saw in the proof of \cref{lem:ChSur}, we can use \cref{lem:FpFn_is_rep} to see that the map $F_p/F_{p-2r} \to F_p/F_{p-r}$ is representable as
    \begin{equation*}
        j^* \colon \Map(\bbS^{-p}_p, K_{[p,p-2r+1]}) \xrightarrow{j \circ } \Map(\bbS^{-p}_p, K_{[p,p-r+1]})
    \end{equation*}
    which is given by postcomposing with restriction and padding. By taking the cofiber of $j$, we get the cofiber of $j^*$. The cofiber is calculated pointwise: for $q \in [p-r,p-2r+1]$, the map $j^*$ is $K_q \to 0$, and for $q \notin [p-r,p-2r+1]$, the map is the identity. We conclude that the cofiber map is of the form:
    \begin{equation*}
        c \colon K_{[p,p-r+1]} \to \cof(j^*) \simeq \S K_{[p-r,p-2r+1]}.
    \end{equation*}
    Hence, the differential is represented by post-composing:
    \begin{equation*}
        c^* \colon \Map(\bbS^{-p}_p, K_{[p,p-r+1]}) \xrightarrow{c \ \circ } \Map(\bbS^{-p}_p, \S K_{[p-r,p-2r+1]}).
    \end{equation*}
    Note that $\bbS^{-p}_p \in \Ch_{(\infty,p]}(\Sp)$ and the other two chains are elements in $\Ch_{[p,-\infty)}(\Sp)$. So we can use \cref{lem:shift_map} on $c$ and find that $c^*$ is equivalent to post-composition with 
    \begin{equation*}
        K_{[p,p-r+1]} \to T^{-1} K_{[p-r,p-2r+1]}
    \end{equation*}
    which is the visualization in the proposition.
    Note that the diagram (which has $K_{p-r}$ as a target) is a visualization of $\partial_r(x) \in \pg{E}{1}{n-1}{p-r}$, and the "non-truncated" version (with target $K_{[p-r,p-2r+1]}$) is the evidence that we can lift $\partial_r(x)$ to $\pg{\cE}{r}{n-1}{p-r}$.
    
    Let $K = \cA F$, we get a cofiber sequence of chain maps:
    \begin{equation*}
        \xymatrix
        {
            0   \ar[r]\ar[d]    & \dots \ar[r]\ar[d]& 0 \ar[r]\ar[d]         &  K_{p-r} \ar[r]\ar[d] & \dots \ar[r]\ar[d]& K_{p-2r+1}\ar[d] \\
            K_p \ar[r]\ar[d]    & \dots \ar[r]\ar[d]& K_{p-r+1}\ar[r]\ar[d] &  K_{p-r} \ar[r]\ar[d] & \dots \ar[r]\ar[d]& K_{p-2r+1}\ar[d] \\
            K_p \ar[r]          & \dots \ar[r]      & K_{p-r+1} \ar[r]      &  0 \ar[r]            & \dots \ar[r]     & 0
        }
    \end{equation*}
    Colimits in the functor category of chains are computed pointwise, hence the connecting morphism is given by
    \begin{equation}\label{eq:diff_sus}
        \xymatrix
        {
            K_p   \ar[r]\ar[d]  & \dots \ar[r]\ar[d]    & K_{p-r+1} \ar[r]\ar[d]    & 0 \ar[r]\ar[d]            & \dots \ar[r]\ar[d] & 0\ar[d] \\
            0 \ar[r]            & \dots \ar[r]          & 0 \ar[r]                  &  \Sigma K_{p-r} \ar[r]    & \dots \ar[r]       &  \Sigma K_{p-2r+1}
        }
    \end{equation}
    We denote the restriction of chains:
    \begin{align*}
        K_{[p,p-r+1]} &:= K_p \to \dots \to K_{p-r+1} \to \overbrace{0 \to \dots \to 0}^{r-1}  \\
        K_{[p-r,p-2r+1]} &:= \overbrace{0 \to \dots \to 0}^{r-1} \to K_{p-r} \to \dots \to K_{p-2r+1}
    \end{align*}
    as objects in $\Ch_{[p,p-2r+2]}(\Sp)$.
    The diagram \cref{eq:diff_sus} is equivalent to the diagram:
    \begin{equation*}
        \xymatrix
        {
            K_{[p,p-r+1]} \ar[r]\ar[d]  & 0 \ar[d] \\
            0 \ar[r]                    & \Sigma K_{[p-r,p-2r+1]}
        }
    \end{equation*}
    Using the stability of $\Sp$, we get that $\Ch_{[p,p-2r+2]}(\Sp)$ is stable, and the last diagram is equivalent to the map of chains: $K_{[p,p-r+1]} \to K_{[p-r,p-2r+1]}$ which can be drawn as:
    \begin{equation*}
        \xymatrix
        {
            K_p   \ar[r]\ar[d]  & \dots \ar[r]\ar[d]    & K_{p-r+1} \ar[r]\ar[d]    & \dots \ar[r]\ar[d]& 0\ar[d] \\
            0 \ar[r]            & \dots \ar[r]          & K_{p-r} \ar[r]            & \dots \ar[r]       & K_{p-2r+1}
        }
    \end{equation*}
    Hence, pre-composition with a lift of $x$ to a map $x\p:\mathbb{S}^{n-p}_p \to K_{[p,p-r+1]}$ gives the diagram:
    \begin{equation}\label{eq:lift_of_dr}
        \xymatrix
        {   \mathbb{S}^{n-p} \ar[r]\ar[d]_x & \dots \ar[r]\ar[d] &  0 \ar[r]\ar[d] & \dots \ar[r]\ar[d] & 0 \ar[d]\\
            K_p \ar[r]\ar[d] & \dots \ar[r]\ar[d] &  K_{p-r+1} \ar[r]\ar[d] & \dots \ar[r]\ar[d] & 0 \ar[d]\\
            0 \ar[r] & \dots \ar[r] & K_{p-r} \ar[r] & \dots \ar[r] & K_{p-2r+1}
        }
    \end{equation}
    and the truncated version:
    \begin{equation*}
        \xymatrix
        {   \mathbb{S}^{n-p} \ar[r]\ar[d] & \dots \ar[r]\ar[d] &  0 \ar[d]\\
            K_p \ar[r]\ar[d] & \dots \ar[r]\ar[d] &  K_{p-r+1} \ar[d] \\
            0 \ar[r] & \dots \ar[r] & K_{p-r} 
        }
    \end{equation*}
    The last diagram is $\partial_r(x) \in \pg{E}{1}{n-1}{p-r}$, and the previous diagram \cref{eq:lift_of_dr} is the evidence that we can lift $\partial_r(x)$ to $\pg{\cE}{r}{n-1}{p-r}$.
\end{proof}